\section{Crossing Number and Writhe}
\subsection{ Crossing Number}
For a diagram $D$, let $c(D)$ be its number of crossings.
The crossing number of a knot $K$ is
\begin{equation}
    c(K) = \min_{D\sim K} c(D),
\end{equation}
where $D \sim K $means that $D$ represents the knot $K$.
Thus,
\[
c(\text{unknot}) =0 ,\quad c(\text{trefoil}) = 3.
\]



\subsection{Signed Crossings}
For an oriented diagram, each crossing is assigned a sign
$\varepsilon_i \in \{+1,-1\}$. The writhe of the diagram is

\begin{equation}
    w(D) = \sum_{i=1}^{c(D)} \varepsilon = n_+(D) - n_-(D)
    \label{eq:three}
\end{equation}

where $n_+(D)$ and $n_-(D)$ count positive and negative crossings. For example, if the signs are $+1, +1, -1, +1$ then
\[
w(D) = 1+1 -1 +1 = 2
\]


\begin{table*}[t]
    \centering
    \caption{ Selected knots and several elementary invariants}
    \label{tab:inv}
    \begin{tabular}{|c|c|c|c|c|c|}
    \hline
     \multirow{2}{*}{\bfseries Family} & \multirow{2}{*}{\bfseries Knot} & \multicolumn{3}{c|}{\bfseries Selected invariants} & \multirow{2}{*}{\bfseries Chiral?} \\
     \cline{3-5}
     & & $c(K)$ & $\det(K)$ & $\Delta_K(t)$ & \\
     \hline
    Basic & Unknot & 0 & 1 & 1 & No \\
    \hline
    \multirow{2}{*}{Torus knots} & Trefoil $3_1$ &3 & 3 & $t^{-1} -1  + t$ & Yes \\
    \cline{2-6}
    & Cinquefoil $5_1$ & 5 & 5 & $t^{-2} - t ^{-1} + 1 - t + t^2$ & Yes\\
    \hline
    Twist knot & Figure-eight  $4_1$ & 4 & 5& $-t^{-1} + 3 - t$ & No \\
    \hline
    \end{tabular}

\end{table*}

Writhe is useful in calculations, but it is not itself a
knot invariant: it can change under a Reidemeister move
of type I.